\section{Weighted measure addition}
\label{sec:big-weighted-limits}

Given $u,v\in\mathcal E_\varphi$ and $\lambda>0$, we look for $a\in\mathcal E_\varphi$ with $\mu_a=\mathrm{e}^{\lambda(a-u)}\mu_u+\mathrm{e}^{\lambda(a-v)}\mu_v$. This weighted addition is used in \cref{sec:big-model-general-measures} to enlarge the weighted measure of a supersolution. The measure $\mathrm{e}^{-\lambda u}\mu_u+\mathrm{e}^{-\lambda v}\mu_v$ can have infinite mass, so we approximate the inputs while keeping a common lower bound in $\mathcal E_\varphi$ for the approximate solutions. This bound gives convergence in capacity, even when the approximate solutions have different ordinary types, and localizes the limiting equation away from deep sublevel sets. Exponentially weighted measures are defined off the pole sets and assigned zero mass there.

\Needspace{13\baselineskip}
\begin{lemma}\label{lem:big-weighted-common-bound}
Let $a_j\in\PSH(X,\theta)$ have a common upper bound and converge in $L^1(X)$ to $a\in\PSH(X,\theta)$ with $a\preceq\varphi$. Suppose there is $h\in\mathcal E_\varphi$ such that $h\leq a_j$ for every $j$. If, for some $\lambda>0$,
\[
 \int_X|a_j-a|\mathrm{e}^{\lambda a}\,\mu_{a_j}\longrightarrow0,
\]
then $a\in\mathcal E_\varphi$ and $a_j\to a$ in capacity.
\end{lemma}
\begin{proof}
By \cref{lem:qpsh-almost-everywhere-order}, the lower bound passes to the limit, so $h\leq a\preceq\varphi$, and upward stability gives $a\in\mathcal E_\varphi$. For the convergence in capacity we follow the proof of \cref{thm:big-relative-self-measure-convergence}, with $a_j,a$ in place of $u_j,u$, the weights $\min(1,|p_{j,N}-a|)\mathrm{e}^{\lambda a}$ in place of $\min(1,|p_{j,N}-u|)$, and the fixed lower bound $h$ in place of $p_1$. Partitioning contact sets gives, along a subsequence with integrals at most $2^{-j}$ and almost-everywhere convergence,
\[
 \int_X\min(1,|p_{j,N}-a|)\mathrm{e}^{\lambda a}\,\mu_{p_{j,N}}
 \leq\sum_{k=j}^N\int_X|a_k-a|\mathrm{e}^{\lambda a}\,\mu_{a_k}
 \leq2^{1-j}
\]
for the finite rooftops $p_{j,N}=P_\theta(a_j,\ldots,a_N)$, which are not identically $-\infty$ because $h$ is a competitor. The first-moment criterion is not needed: the decreasing limits $p_j$ satisfy $h\leq p_j\leq a$, so the $p_j$ and their increasing regularized limit $p$ lie in $\mathcal E_\varphi$, although the finite rooftops need not. The weights converge in capacity: discard deep sublevel sets of $h$ and $a$, and use continuity on the remaining finite ranges. Since the exponential is positive off a pluripolar set, which $\mu_p$ does not charge, the limiting identity $\int_X\min(1,|p-a|)\mathrm{e}^{\lambda a}\,\mu_p=0$ again gives $\mu_p(\{p<a\})=0$, and the rest of that proof applies.
\end{proof}

The next lemma treats exponent one; the general exponent in \cref{thm:big-weighted-addition} follows by rescaling the background and the potentials.

\Needspace{20\baselineskip}
\begin{lemma}\label{lem:big-truncated-addition-limit}
Let $u,v,\sigma\in\PSH(X,\theta)$ be nonpositive, with $u,v\preceq\varphi$, and let $t_j\uparrow\infty$. Each of the sequences $u_j$, $v_j$ is either constant, equal to $u$, respectively $v$, for every $j$, or given by the truncations
\[
 u_j=\max(u,\sigma-t_j),\qquad
 v_j=\max(v,\sigma-t_j).
\]
Fix $h\in\mathcal E_\varphi$ and let $a_j\in\PSH(X,\theta)$ satisfy
\[
 \mu_{a_j}=\mathrm{e}^{a_j-u_j}\mu_{u_j}+\mathrm{e}^{a_j-v_j}\mu_{v_j},
 \qquad h\leq a_j\leq\min(u_j,v_j).
\]
If $a_j\to a$ in $L^1(X)$, then $a\in\mathcal E_\varphi$, the convergence holds in capacity, and
\[
 a\leq\min(u,v),\qquad
 \mu_a=\mathrm{e}^{a-u}\mu_u+\mathrm{e}^{a-v}\mu_v.
\]
\end{lemma}
\begin{proof}
Since $u_j$ and $v_j$ decrease to $u$ and $v$, the bounds give $h\leq a\leq\min(u,v)\preceq\varphi$, so $a\in\mathcal E_\varphi$. After discarding finitely many indices we may assume $t_j\geq0$, so all potentials $a_j,a$ are nonpositive. Let $E_j$ be the intersection of the sets $\{u>\sigma-t_j\}$ and $\{v>\sigma-t_j\}$, taken only for those inputs that are truncated; a constant sequence imposes no condition. On $E_j$, locality gives $\mu_{u_j}=\mu_u$ and $\mu_{v_j}=\mu_v$. Since $a,a_j\leq u_j,v_j$, we have
\[
 \mathrm{e}^a\mu_{a_j}
 \leq \mathrm{e}^{(a+a_j)/2}(\mu_{u_j}+\mu_{v_j}).
\]
If $x\notin E_j$, then one truncated input and its truncation are at most $-t_j$ at $x$, so $a_j(x),a(x)\leq-t_j$. Moreover $|a_j-a|\mathrm{e}^{(a+a_j)/4}$ is bounded by an absolute constant, since $a_j,a\leq0$. Thus
\begin{equation}\label{eq:big-addition-weighted-error}
 \int_X|a_j-a|\mathrm{e}^a\,\mu_{a_j}
 \leq\int_X|a_j-a|\mathrm{e}^{(a+a_j)/2}(\mu_u+\mu_v)
       +C \mathrm{e}^{-t_j/2}.
\end{equation}
The last term uses the uniform bound for the masses of the $u_j$ and $v_j$ given by \cref{prop:bmp-finite}. The fixed measure $\mu_u+\mu_v$ is finite by bounded mass. \Cref{lem:fixed-npp-integrals}, applied to $\mathrm{e}^{a_j}$ and $\mathrm{e}^a$, shows that $a_j\to a$ in measure for $\mu_u+\mu_v$, since $a$ is finite almost everywhere for this non-pluripolar measure. The integrand in the first term of \eqref{eq:big-addition-weighted-error} is uniformly bounded and tends to zero in measure, so its integral tends to zero. \Cref{lem:big-weighted-common-bound} now gives convergence in capacity. When a subsequence is taken, the corresponding truncation levels $t_j$ are kept.

To identify the equation, fix $T>0$ and choose a continuous nondecreasing function $g_T\colon[-\infty,0]\to[0,1]$ that vanishes on $[-\infty,-2T]$ and equals one on $[-T,0]$, such that $g_T\uparrow1$ at finite points as $T\to\infty$. Put $\chi_T=g_T(h)$. Since $h\leq a\leq u,v$, both $u$ and $v$ exceed $-2T$ on $\{\chi_T>0\}$. Therefore, once $t_j>2T$, every truncated input agrees there with its limit, and
\[
 \chi_T\mu_{a_j}
 =\chi_T\mathrm{e}^{a_j}(\mathrm{e}^{-u}\mu_u+\mathrm{e}^{-v}\mu_v).
\]
\Needspace{7\baselineskip}
Choose a Hermitian form $\Omega\geq\theta$. On the left-hand side, replace $a_j$ by $\max(a_j,-2T)$, which does not change $\chi_T\mu_{a_j}$ because $\chi_T$ vanishes where $a_j\leq-2T$, and expand the bounded product as
\[
 (\theta+\ddc\max(a_j,-2T))^n
 =\sum_{k=0}^n\binom nk
   (\Omega+\ddc\max(a_j,-2T))^k\wedge(\theta-\Omega)^{n-k}.
\]
Bounded capacity continuity with the fixed quasi-continuous weight $\chi_T$ therefore gives convergence to $\chi_T\mu_a$, after testing against a continuous function. The uncut product may be signed, but its terms have uniformly capacity-controlled total variations.

On the right-hand side, $\chi_T \mathrm{e}^{-u}$ and $\chi_T \mathrm{e}^{-v}$ are bounded by $\mathrm{e}^{2T}$. Fixed-measure exponential convergence therefore gives convergence in total variation to
$\chi_T\mathrm{e}^a(\mathrm{e}^{-u}\mu_u+\mathrm{e}^{-v}\mu_v)$. Let $T\to\infty$. The limiting right-hand side is at most $\mu_u+\mu_v$, because $a\leq u,v$. Both sides ignore the pole set of $h$, so the cutoff identities prove the global equation.
\end{proof}
